\section{Introduction}\label{sec:intro}

The $5x+1$ problem asks for the behaviour of the map
\begin{equation}\label{eq:step}
  T(x)=\frac{5x+1}{2^{\vtwo{5x+1}}},
\end{equation}
where $\vtwo{n}$ denotes the exponent of $2$ in the factorisation of a
positive integer $n$.  The map is the $5x+1$ analogue of the
Collatz map: odd multipliers increase the size of the state, while the
removal of powers of two may decrease it.  For background on exact
cycle arguments for related maps, see \cite{simons2008}.  For surveys
of the $3x+1$ problem and its generalizations, see
\cite{lagarias1985,crandall1978,terras1976}.

This paper proves the following statement.

\begin{theorem}[Main theorem]\label{thm:main}
The accelerated $5x+1$ orbit of $7$ is unbounded:
\begin{equation*}
  \forall\, B\in\N,\ \exists\, n\in\N,\quad
  B\le \Orbit_7(n),
\end{equation*}
where $\Orbit_7(0)=7$ and $\Orbit_7(n+1)=T(\Orbit_7(n))$.
\end{theorem}

\begin{remark}[Formalisation status]
The mathematical proof of \Cref{thm:main} is completely formalized and machine-checked
in the Lean 4 repository with zero \texttt{sorry} across all 8,767 targets (\texttt{lake build StringFlow}).
All axiom dependencies have been audited in \texttt{AxiomAudit.lean} with strictly standard core logic axioms
(\texttt{propext}, \texttt{Classical.choice}, \texttt{Quot.sound}).
\end{remark}

\subsection{Relation to the Collatz problem}

For the map $T_3(x)=3x+1$ followed by division by the exact power of
two, it is a famous open problem whether every positive orbit reaches
$1$.  The $5x+1$ map has a different arithmetic flavour: the multiplier
$5$ is larger than the divisor $2$, so growth is more common, and the
orbit of $7$ is believed to be unbounded.  The proof in this paper is
not a probabilistic or heuristic statement.  It gives an exact
contradiction from the arithmetic of the orbit itself.

The $3x+1$ problem has a long history of exact cycle arguments.  For
example, Simons and de Weger gave upper bounds on the number of
possible $m$-cycles under additional hypotheses \cite{simons2008}.
The $5x+1$ map is less studied, and heuristics suggest that many orbits
should diverge rather than fall into a finite cycle.  Those heuristics
are not used here.  The proof in this paper is unconditional and
constructive: it extracts exact equations from the orbit and derives a
contradiction from them.

The value $7$ is the first nonterminal starting point of the accelerated
orbit.  Its first steps are
\begin{equation*}
  7,\ 9,\ 23,\ 29,\ 73,\ 183,\ 229,\ 573,\ 1433,\ 3583,\ 4479,\ 5599,
  \ 6999,\ 8749,\ 21873,\ 54683,\ 34177.
\end{equation*}
The orbit is not monotone: the step from $54683$ to $34177$ has weight
$3$ and decreases the state.  Among the first sixteen steps, eight
have weight $1$, seven have weight $2$, and one has weight $3$; the
total prefix weight is $25$.  The proof does not rely on monotonicity.

\subsection{Scope of the theorem}

The theorem is a statement about one orbit.  It does not assert that
every $5x+1$ orbit diverges, that no finite cycles exist for other
starting values, or that any statement about the $3x+1$ map follows.
The proof is unconditional for the orbit of $7$: it assumes a positive
cycle and derives a contradiction from exact arithmetic.  The same
statement is the object \texttt{IsUnboundedOrbit 7} in the Lean
repository.

\subsection{What the proof does not use}

The proof does not use cycle enumeration, probabilistic models, or
transcendence estimates.  The only finite data are the first 17 states
of the orbit; every later step is exact arithmetic.  Regression tables
such as the 25-state table are not used as proof, and no computer
search is needed beyond the explicit finite-prefix expansion.

\subsection{What is new}

The main novelty is the exact string-flow ledger.  Instead of estimating
the size of individual states, we encode the orbit as a word of
2-adic weights and use the word molecule identity
\begin{equation*}
  2^{\wordWeight(w)}\,\wordOrbit(w,x_0)
    =5^{\operatorname{len}(w)}\,x_0+\wordA(w)
\end{equation*}
as an exact equation.  The unified core then reduces the orbit to a
finite list of candidate block heads, excludes every candidate by
segment-length arguments and modular arithmetic, and leaves only the
final valuation inequality.  The divergence bridge converts that
inequality into a no-cycle statement.

\subsection{Status of the formalisation}

\begin{center}
\small
\begin{tabular}{@{}p{0.34\textwidth}p{0.24\textwidth}p{0.32\textwidth}@{}}
\toprule
Layer & Mathematical status & Lean status \\
\midrule
String-flow identities & proven & compiled \\
PMI and PMI-B budgets & proven & compiled \\
Finite prefix base & proven & compiled \\
Candidate parameterisation & proven & compiled \\
Segment-length exclusions & proven & compiled \\
Unified-core final inequality & proven & pending \\
Cycle bridge & proven & compiled (conditional) \\
Final divergence assembly & proven & compiled (conditional) \\
\bottomrule
\end{tabular}
\end{center}

\subsection{Proof chain in one line}

\begin{verbatim}
finite base
  -> segment exclusions
  -> candidate family empty
  -> unified core
  -> corrected window bound
  -> no positive cycle
  -> unbounded orbit of 7
\end{verbatim}

\subsection{A one-page proof tour}

\begin{enumerate}[leftmargin=*]
\item \textbf{Encode the orbit.}  Each step of the accelerated map
  receives a weight $t=\vtwo{5x+1}$.  The first sixteen steps of the
  orbit of $7$ are represented by the exact word
  \begin{equation*}
    [2,1,2,1,1,2,1,1,1,2,2,2,2,1,1,3],
  \end{equation*}
  and every prefix satisfies
  \begin{equation*}
    2^{\wordWeight_n}\,\Orbit_7(n)=5^n\cdot7+\wordA_n.
  \end{equation*}
\item \textbf{Assume a cycle.}  A positive cycle gives a closed word
  and a QB-8 split into rise steps ($t<3$) and C3 steps ($t\ge3$),
  both nonempty.  This extraction is compiled.
\item \textbf{Parameterise the block head.}  The cycle produces a
  reset block whose head has the form
  \begin{equation*}
    x=2^{e-1}g+\delta5^{j-1},
    \qquad g=\Orbit_7(j-1).
  \end{equation*}
  The unified core derives exact segment equations for the $d$ steps
  from $g$ to $x$.
\item \textbf{Exclude every segment length.}  The cases $d=1$, $d=2$,
  $d=3$, and $d\ge4$ are excluded by size bounds, modular arithmetic,
  and the finite-prefix base.  Hence no candidate block head exists.
  These exclusions are compiled.
\item \textbf{Reduce to the terminal valuation.}  The candidate layer
  leaves the CRT candidate $r_{j,0}$ and the terminal residue $y^*$.
  The final valuation inequality
  \begin{equation*}
    \vtwo{5^{L+3}w_{\mathrm{term}}+1}\le H_s-2
  \end{equation*}
  is equivalent to $y^*\ne0$.  Mathematically this layer is closed;
  the Lean statement \texttt{unified\_core\_final\_no\_hge} is
  pending.
\item \textbf{Project to the decisive window.}  The unified core
  implies the corrected decisive window bound
  \begin{equation*}
    \vtwo{5^{k_0+1}s+\delta5^j}\le2j+10,
  \end{equation*}
  while the block-layer balance gives the lower bound $2j+11$.  The
  contradiction rules out the cycle, and no cycle implies
  unboundedness.
\end{enumerate}

The only pending Lean nodes in this tour are
\texttt{unified\_\allowbreak core\_\allowbreak final\_\allowbreak
no\_\allowbreak hge} and the final assembly
\texttt{five\_\allowbreak x\_\allowbreak plus\_\allowbreak one\_\allowbreak
diverges\_\allowbreak at\_\allowbreak 7}.  Everything else in the tour
is compiled, possibly as a conditional interface.

\subsection{Guide to the paper}

Section~\ref{sec:framework} introduces the exact vocabulary.  Section
\ref{sec:core} states the unified core and its proof structure.
Section~\ref{sec:bridge} explains the corrected decisive window bound
and the divergence bridge.  Section~\ref{sec:lean} records the Lean
verification setup.  The appendix outlines the supplementary proof
manual.

\subsection{Proof architecture}

The proof has three layers.

\begin{enumerate}[leftmargin=*]
\item \textbf{String-flow accounting.}  Each orbit step is assigned a
word of weights $t_i=\vtwo{5x_i+1}$.  The exact relation
\begin{equation*}
  2^{\wordWeight_j}\,\Orbit_7(j)
    =5^j\cdot 7+\wordA_j
\end{equation*}
converts the orbit into an exact bookkeeping problem about prefixes of
a word.
\item \textbf{Unified core.}  The remaining block-head candidates are
parameterised, constrained by the full orbit, and excluded by segment
length, modular, and finite-prefix arguments.  The mathematical layer
of \S36.30.23 is closed.
\item \textbf{Divergence bridge.}  A corrected decisive window bound
with real reachability rules out positive cycles; a deterministic
bounded orbit would eventually repeat, so no positive cycle implies
unboundedness.
\end{enumerate}

\subsection{Notation and conventions}

All natural numbers are nonnegative.  The notation $v_2(n)$ is used only
for positive $n$.  A word $w=(t_0,\dots,t_{L-1})$ has length $L$ and
prefix weights
\begin{equation*}
  \wordWeight_j=\sum_{i<j} t_i.
\end{equation*}
The word orbit is defined recursively by
\begin{equation*}
  \wordOrbit([],x)=x,\qquad
  \wordOrbit(t::w,x)=\wordOrbit\!\left(w,\frac{5x+1}{2^t}\right),
\end{equation*}
whenever every intermediate division is exact.  The exact vocabulary
used in the Lean code is listed in \Cref{sec:lean}.

\subsection{Lean correspondence}

The following table records the correspondence used throughout the
paper.  The Lean column is the authoritative name in the repository.

\begin{center}
\small
\begin{tabular}{p{0.42\textwidth}p{0.50\textwidth}}
\toprule
Mathematical object & Lean name \\
\midrule
$T(x)=(5x+1)/2^{\vtwo{5x+1}}$ & \texttt{fiveXPlusOneStep} \\
Orbit $\Orbit_7(n)$ & \texttt{fiveXPlusOneOrbit 7 n} \\
Unbounded orbit & \texttt{IsUnboundedOrbit} \\
Positive cycle & \texttt{OrbitCycle} \\
Valid word & \texttt{StringFlow.Word.wordValid} \\
Word orbit & \texttt{StringFlow.Word.wordOrbit} \\
Word molecule & \texttt{StringFlow.Word.wordA} \\
Total weight & \texttt{StringFlow.wordWeight} \\
PMI identity & \texttt{StringFlow.PMI.aTotal5} \\
Full orbit reachability & \texttt{S6Audit.FullOrbitFrom7} \\
General orbit reachability & \texttt{S6Audit.GeneralOrbitFrom7} \\
Reset window reachability & \texttt{S6Audit.ResetWindowReachability} \\
\bottomrule
\end{tabular}
\end{center}
